\section*{Reproducing the results}

\textit{What is in the repository.} The original runs' accuracy files are in \texttt{results/\allowbreak original/}, and \texttt{analysis/\allowbreak recompute\_original.py} recomputes \texttt{analysis/\allowbreak out/\allowbreak summary\_original.csv} from them alone. The CPU checks (\texttt{analysis/\allowbreak checks.py}: raw pixels, the rerun of the saved model, and the detector re-creation) need the original runs' saved CelebA-HQ samples, which are not in the repository, so their committed outputs in \texttt{analysis/\allowbreak out/} cannot be regenerated from it. The original model's weights (\texttt{model.pth}), which the rerun and the study's original-model arm need, are attached to the repository's release, with SHA-256 \texttt{\N{orig_sha256}}. The study's code is in \texttt{study/}, and its aggregate results from the run of November 11, 2026, are in \texttt{results/\allowbreak study/}: \texttt{summary.csv}, the seven analysis files, each model's fixed thresholds, the attention checks, the training logs, \texttt{run\_log.json} (settings, environment, and hashes), and \texttt{run.log}. \texttt{near\_copies.csv} is written by \texttt{study/\allowbreak check\_near\_copies.py} from the per-probe results, which are not in the repository, so it cannot be regenerated from it either. The run took \N{wall_s}~s (\N{wall_h}~h \N{wall_m}~min) for \N{units} units: \N{wall_s1}~s for alignment and the pretrained models on the CPU, then \N{wall_s2}~s for the rest, including \N{train_s}~s of training (\N{train_s_lo} to \N{train_s_hi}~s per run). It used Python \N{py}, PyTorch \N{torch_v}~\cite{paszke2019pytorch} with CUDA \N{cuda}, OpenCV \N{cv2_v}~\cite{bradski2000opencv}, NumPy \N{np_v}~\cite{harris2020numpy}, pandas \N{pd_v}~\cite{mckinney2010pandas}, and Pillow \N{pil_v}~\cite{pillow}. The original runs used DeepFace \N{deepface_v}~\cite{deepface} with TensorFlow \N{tf_v}~\cite{abadi2016tensorflow}, rembg \N{rembg_v}~\cite{rembg}, and ONNX Runtime \N{ort_v}~\cite{onnxruntime} on Google Colab; their PyTorch version was not recorded. The bootstrap analysis (\texttt{study/\allowbreak analyze.py}) ran on Python \N{an_py} with NumPy \N{an_np} and pandas \N{an_pd}. The study's training runs are seeded, but cuDNN's deterministic mode was off, so a rerun may not reproduce the committed results exactly, and the repository records no such comparison.

\textit{What is not.} No CelebA-HQ image, identity label, or file name is in the repository. The study's per-probe results, its split, and its aligned crops carry CelebA file names and identity numbers and are left out, so \texttt{study/\allowbreak analyze.py}, which reads the per-probe results, cannot regenerate the committed analysis files from the repository alone. The \N{n_runs} trained study models and the image map are not included. \texttt{study/\allowbreak README.md} gives the sources of the data, the label file, and the weights, with the image set's fingerprint and the files' SHA-256 hashes, and how to rebuild the image map with \texttt{match\_images.py} and check it with \texttt{check\_labels.py} before a full run. \texttt{run\_study.py -{}-smoke} runs a small synthetic version of the pipeline (generated images, mock pretrained models, and a randomly initialized original model) on a CPU, with no data or weights.

\textit{Data and model terms.} CelebA~\cite{liu2015celeba}, from which CelebA-HQ~\cite{karras2018pggan} is built, is available for non-commercial research only and may not be further distributed; CelebA-HQ, whose creators released the files that rebuild it from CelebA under CC BY-NC 4.0, is used here on the same terms, and CelebA-Dialog's data~\cite{jiang2021talk} carry the same terms. The Kaggle copy the study used~\cite{kagglecelebahq} carries a license label from its uploader, which does not change CelebA's terms. The original model's weights were trained on CelebA-HQ images and are released for non-commercial research only. InsightFace's pretrained models, including \texttt{w600k\_r50}, are for non-commercial research only~\cite{insightface}. SFace and YuNet are distributed through the OpenCV Zoo~\cite{opencvzoo}, SFace under the Apache License and YuNet under the MIT License.

\textit{How the paper's numbers are made.} Every result number and table here is generated by \texttt{paper/\allowbreak make\_numbers.py} from the committed files above, and every figure is drawn from them by \texttt{paper/\allowbreak make\_figures.py} with Matplotlib~\cite{hunter2007matplotlib}; no figure uses an image. Fig.~\ref{fig:occluders} draws the occluders from the study's settings, with the five landmark positions of the alignment template typed in from OpenCV's \texttt{alignCrop} source. The one exception typed into \texttt{make\_numbers.py} is the estimate that about 1 in \N{orig_overlap_share} of the study's gallery frames was among the original model's training images; a few statements of which condition, seed, or sample a number belongs to are written into the text and were checked against the files. The study's settings come from \texttt{study/\allowbreak config.json} and \texttt{run\_log.json}; the original model's training settings (\N{orig_train_images} images, \N{orig_epochs} epochs, batch size \N{orig_batch}, learning rate \N{orig_lr}, margin \N{orig_margin}, and squares of \N{orig_occ_lo} to \N{orig_occ_hi}\%) from \texttt{frfpv1/\allowbreak config.py}, which keeps the values of the original notebook (not in the repository); and the parameter counts from \texttt{frfpv1/\allowbreak models.py}.
